\documentclass[11pt]{article}
\usepackage{amsmath,amssymb,amsthm}
\usepackage{hyperref}

\newtheorem{theorem}{Theorem}[section]
\newtheorem{lemma}[theorem]{Lemma}
\newtheorem{definition}[theorem]{Definition}
\newtheorem{corollary}[theorem]{Corollary}

\begin{document}

\title{On a Polynomial Identity in Hyperfields}
\author{}
\date{}
\maketitle

\begin{abstract}
We answer the question whether, in an arbitrary hyperfield $\mathbb{H}$, the condition
$0 \in x \boxplus y \boxplus z$ implies
$0 \in x^3 \boxplus y^3 \boxplus z^3 \boxplus -xyz \boxplus -xyz \boxplus -xyz$.
We prove that the implication is true. The proof uses Marshall's representation theorem,
which states that every hyperfield is a quotient of a commutative ring by a subgroup of its
group of units. In such a quotient, the classical polynomial identity
$x^3+y^3+z^3-3xyz = (x+y+z)(x^2+y^2+z^2-xy-yz-zx)$ lifts to the hyperfield setting.
\end{abstract}

\section{Introduction and Statement of the Result}

Let $\mathbb{H}$ be a hyperfield with multivalued addition $\boxplus$ and multiplication $\cdot$.
Recall that a hyperfield is a tuple $(\mathbb{H},0,1,\cdot,\boxplus)$ such that
$(\mathbb{H},0,1,\cdot)$ is a commutative monoid with $0$ absorbing,
$(\mathbb{H},\boxplus,0)$ is a canonical hypergroup (commutative, associative, with identity
$0$, unique inverses, and reversibility), and multiplication distributes over hyperaddition:
\[
a \cdot (b \boxplus c) = (a\cdot b) \boxplus (a\cdot c)
\]
for all $a,b,c \in \mathbb{H}$, where the right-hand side is understood as the union
$\bigcup_{d \in b \boxplus c} a\cdot d$.

The question under consideration is the following.

\begin{quote}
\itshape
Let $\mathbb{H}$ be a hyperfield, and let $x,y,z \in \mathbb{H}$.
Suppose that $0 \in x \boxplus y \boxplus z$.
Does it then follow that
\[
0 \in x^3 \boxplus y^3 \boxplus z^3 \boxplus -xyz \boxplus -xyz \boxplus -xyz \text{?}
\]
\end{quote}

We answer this question affirmatively.

\begin{theorem}\label{thm:main}
Let $\mathbb{H}$ be a hyperfield. For any $x,y,z \in \mathbb{H}$, if
$0 \in x \boxplus y \boxplus z$, then
\[
0 \in x^3 \boxplus y^3 \boxplus z^3 \boxplus -xyz \boxplus -xyz \boxplus -xyz.
\]
\end{theorem}

\section{Preliminaries: Quotient Hyperfields}

We use the following standard representation theorem for hyperfields.

\begin{theorem}[Marshall's Representation Theorem]\label{thm:marshall}
Every hyperfield $\mathbb{H}$ is isomorphic to a quotient hyperfield $R/G$, where
$R$ is a commutative ring and $G$ is a subgroup of the group of units $R^\times$ of $R$.
The quotient hyperfield $R/G$ is defined as follows.
\begin{itemize}
\item The elements of $R/G$ are the orbit sets $[a] := \{a g : g \in G\}$ for $a \in R$,
      together with the zero class $[0]=\{0\}$.
\item Multiplication is defined by $[a]\cdot [b] := [ab]$.
\item Hyperaddition is defined by
      \[
      [a] \boxplus [b] := \{ [a' + b'] : a' \in [a],\ b' \in [b] \}.
      \]
\end{itemize}
The additive inverse of $[a]$ is $[-a]$.
\end{theorem}

For a proof and further discussion of this theorem, see Marshall's work on hyperfields
\cite{Marshall}. The theorem is also mentioned in the context of multigroups in
\cite{Viro} and \cite{ViroTropical}. In particular, every hyperfield can be realized as
such a quotient, so it suffices to prove the desired implication in the quotient model.

\begin{remark}
In the quotient model $R/G$, the condition $0 \in [a] \boxplus [b] \boxplus [c]$ means
that there exist elements $a' \in [a]$, $b' \in [b]$, and $c' \in [c]$ such that
$a' + b' + c' = 0$ in the ring $R$.
\end{remark}

\section{Proof of the Main Theorem}

\begin{proof}[Proof of Theorem~\ref{thm:main}]
Let $\mathbb{H}$ be a hyperfield. By Theorem~\ref{thm:marshall}, we may assume
$\mathbb{H} = R/G$ for some commutative ring $R$ and a subgroup $G \subseteq R^\times$.

Let $x,y,z \in \mathbb{H}$. Choose representatives $a,b,c \in R$ such that
$x = [a]$, $y = [b]$, and $z = [c]$.

Assume that $0 \in x \boxplus y \boxplus z = [a] \boxplus [b] \boxplus [c]$.
By the definition of hyperaddition in the quotient, there exist $a' \in [a]$,
$b' \in [b]$, and $c' \in [c]$ such that
\[
a' + b' + c' = 0 \quad \text{in } R.
\]
Note that $a' \in [a]$ means $a' = a g_a$ for some $g_a \in G$, and similarly
$b' = b g_b$, $c' = c g_c$ for some $g_b, g_c \in G$.

Now consider the classical polynomial identity in the commutative ring $R$:
\[
a'^3 + b'^3 + c'^3 - 3a'b'c' = (a' + b' + c')(a'^2 + b'^2 + c'^2 - a'b' - b'c' - c'a').
\]
Since $a' + b' + c' = 0$, the right-hand side is $0$. Hence
\[
a'^3 + b'^3 + c'^3 - 3a'b'c' = 0.
\]

We now pass to the quotient hyperfield $R/G$. The class of $a'^3$ is
$[a'^3] = [a]^3 = x^3$, because $a' \in [a]$ and multiplication is well-defined on
quotient classes. Similarly, $[b'^3] = y^3$ and $[c'^3] = z^3$.
Also, $[a'b'c'] = [a][b][c] = xyz$. Therefore, the class of the element
$a'^3 + b'^3 + c'^3 - 3a'b'c'$ belongs to the hyperaddition
\[
x^3 \boxplus y^3 \boxplus z^3 \boxplus (-xyz) \boxplus (-xyz) \boxplus (-xyz).
\]
Indeed, the element $a'^3$ is a representative of $x^3$, $b'^3$ of $y^3$,
$c'^3$ of $z^3$, and $-a'b'c'$ is a representative of $-xyz$. Since the sum of these
representatives is $0$ in $R$, the class $[0] = 0$ lies in the hyperaddition.
Thus
\[
0 \in x^3 \boxplus y^3 \boxplus z^3 \boxplus -xyz \boxplus -xyz \boxplus -xyz.
\]
This completes the proof.
\end{proof}

\section{Conclusion}

We have shown that the implication holds in every hyperfield. The proof relies on
Marshall's representation theorem, which allows us to work in a quotient of a
commutative ring where the classical polynomial identity is valid. The multivalued
nature of the addition in the hyperfield is precisely captured by the fact that in the
quotient model we may choose representatives independently for each occurrence of a
class.

\begin{thebibliography}{9}

\bibitem{Marshall}
M. Marshall,
\emph{Real reduced hyperfields},
J. Algebra \textbf{102} (1986), 1--17.

\bibitem{Viro}
O. Viro,
\emph{Hyperfields for Tropical Geometry I. Hyperfields and dequantization},
arXiv:1006.3034.

\bibitem{ViroTropical}
O. Viro,
\emph{Hyperfields and dequantization},
Stony Brook preprint, available at
\url{https://www.math.stonybrook.edu/~oleg/VersoTeX/Show/Hyperfields/hf.html}.

\end{thebibliography}

\end{document}